\section{Conclusion}

This paper has presented a scalable, fully unsupervised deep learning framework designed to bridge the verification gap in sustainable commercial building operations. By addressing the critical challenge of high-cost hardware sub-metering, our approach empowers facility managers to audit lighting efficiency using only readily available, aggregate hourly electricity meter data.

Our methodology combines a 1D Convolutional Autoencoder (1D-CAE) for robust operational pattern clustering with an adapted N-BEATS architecture for single-channel signal disaggregation. The 1D-CAE successfully extracts a compact latent manifold from daily load profiles, enabling the identification of distinct building control regimes (e.g., scheduled versus reactive use) with a Silhouette Score of 0.3728. Building upon this contextual understanding, the N-BEATS model decomposes the aggregate signal into a continuous base load (primarily HVAC) and discrete, high-frequency pulses (primarily lighting).

Extensive evaluation on the public ASHRAE Great Energy Predictor III dataset—encompassing nearly 500,000 daily profiles from 1,448 buildings across 16 climate zones—demonstrates the framework's superior performance over classical decomposition methods and supervised baselines (MSE of 0.0628). Furthermore, our experiments confirm the model's resilience to sensor noise and its critical ability to generalize across distinct geographic regions without site-specific retraining.

The practical case study on an educational facility illustrates the immediate real-world utility of this approach. By isolating the lighting load, the framework provides granular, actionable insights into operational inefficiencies, allowing building operators to optimize automated lighting schedules and identify faulty sensors. This data-driven visibility is essential for ensuring that the designed energy efficiency of commercial buildings is actually achieved in practice.

Future work will focus on extending the unsupervised N-BEATS architecture to disaggregate additional end-use categories, such as plug loads and specific HVAC sub-components (e.g., chillers versus fans). Additionally, integrating real-time weather data into the disaggregation process could further refine the separation of temperature-dependent base loads from occupancy-driven consumption. Ultimately, the transparent and reproducible methodology presented here provides a powerful tool for the applied data science community to advance the goals of sustainable architecture and energy conservation.

\section*{Acknowledgments}
This study utilizes the ASHRAE Great Energy Predictor III dataset provided via Kaggle. We thank the organizers for making this data publicly available for research.

\begin{thebibliography}{10}

\bibitem{ashrae_2020}
Miller, C., et al.: The Building Data Genome Project 2, energy meter data from the ASHRAE Great Energy Predictor III competition. Scientific Data, 7(1), 1-15 (2020).

\bibitem{performance_gap_1}
De Wilde, P.: The gap between predicted and measured energy performance of buildings: A framework for investigation. Automation in Construction, 41, 40-49 (2014).

\bibitem{performance_gap_2}
Menezes, A. C., et al.: Predicted vs. actual energy performance of non-domestic buildings: Using post-occupancy evaluation data to reduce the performance gap. Energy and Buildings, 47, 355-364 (2012).

\bibitem{performance_gap_3}
Bordass, D., et al.: The performance gap: causes and solutions. Building Research \& Information, 32(5), 414-429 (2004).

\bibitem{eia_lighting}
U.S. Energy Information Administration (EIA): Commercial Buildings Energy Consumption Survey (CBECS) Data. \url{https://www.eia.gov/consumption/commercial/} (2018).

\bibitem{hart_1992}
Hart, G. W.: Nonintrusive appliance load monitoring. Proceedings of the IEEE, 80(12), 1870-1891 (1992).

\bibitem{nilm_survey_1}
Zoha, A., et al.: Non-intrusive load monitoring approaches for disaggregated energy sensing: A survey. Sensors, 12(12), 16838-16866 (2012).

\bibitem{nilm_survey_2}
Faustine, A., et al.: A survey on non-intrusive load monitoring methodies and techniques for energy disaggregation problem. IEEE Access, 5, 210142-210168 (2017).

\bibitem{stl_1990}
Cleveland, R. B., et al.: STL: A seasonal-trend decomposition procedure based on loess. Journal of Official Statistics, 6(1), 3-73 (1990).

\bibitem{fhmm_2012}
Kolter, J. Z., Jaakkola, T.: Approximate inference in additive factorial HMMs with application to energy disaggregation. In: Artificial Intelligence and Statistics, pp. 1472-1482 (2012).

\bibitem{kelly_2015}
Kelly, J., Knottenbelt, W.: Neural NILM: Deep neural networks applied to energy disaggregation. In: Proceedings of the 2nd ACM International Conference on Embedded Systems for Energy-Efficient Built Environments, pp. 55-64 (2015).

\bibitem{zhang_2018}
Zhang, C., et al.: Sequence-to-point learning with neural networks for non-intrusive load monitoring. In: Proceedings of the AAAI Conference on Artificial Intelligence, 32(1) (2018).

\bibitem{bert4nilm}
Yue, Z., et al.: BERT4NILM: A bidirectional transformer model for non-intrusive load monitoring. In: Proceedings of the 7th ACM International Conference on Systems for Energy-Efficient Buildings, Cities, and Transportation, pp. 289-298 (2020).

\bibitem{cae_2021}
Thill, M., et al.: Temporal convolutional autoencoder for unsupervised anomaly detection in time series. Applied Soft Computing, 112, 107751 (2021).

\bibitem{nbeats_2020}
Oreshkin, B. N., et al.: N-BEATS: Neural basis expansion analysis for interpretable time series forecasting. In: International Conference on Learning Representations (ICLR) (2020).

\bibitem{redd_2011}
Kolter, J. Z., Johnson, M. J.: REDD: A public data set for energy disaggregation research. In: Workshop on Data Mining Applications in Sustainability (SIGKDD), 25(59-62), 114-119 (2011).

\bibitem{ukdale_2015}
Kelly, J., Knottenbelt, W.: The UK-DALE dataset, domestic appliance-level electricity demand and whole-house demand from five UK homes. Scientific Data, 2(1), 1-14 (2015).

\end{thebibliography}
